% ---
% id: receiver/rem:tolerancedomain
% kind: rem
% title: "The domain of the receiver tolerances"
% summary: >-
%   Every quantity a receiver tolerance bounds, decides or reduces lies in
%   [0, 1], so at a tolerance of one or more its test is vacuous or degenerate.
%   The declared domain [0, 1) makes every error term the contract states hold
%   on every value it admits.
% prominence: standard
% status: accepted
% uses: [def:residuals, def:tolmap, thm:emitted-bound, lem:exactbasic, rem:exacteig, def:pmcex, thm:multiclick-bound]
% mentions: [rem:ranktwosided, rem:cdirection, receiver/detector-imperfections/prop:dark-slack-bound, req:tolerance-domain, def:setting-family]
% assumes: []
% refs: []
% provenance: newly-derived
% lean:
% ---

\begin{remark}
\label{rem:tolerancedomain}
Each of $\varepsilon_1,\dots,\varepsilon_6$, $\varepsilon_{\mathrm{gap}}$ and
$\varepsilon_{\mathrm{eig}}$ is declared in $[0,1)$ and validated at input
(Requirement~\ref{req:tolerance-domain}). A declared number needs a declared
domain for the same reason a declared parameter does, and here the domain is
not a convention: \emph{a \tref[receiver tolerances]{receiver tolerance} is a
slack on a normalized quantity, never a scale.} Every quantity the eight bound,
decide or reduce already lies in $[0,1]$: eigenvalues of operators between $0$
and $\id$ (\eqref{eq:res2}, \eqref{eq:eps2rank},
Definition~\ref{def:tolmap}\,(i),\,(v)), the spectral-norm defects
\eqref{eq:res1}, \eqref{eq:res3}, \eqref{eq:res4} and \eqref{eq:rhodef} of
operators bounded by $\id$, entries and column defects of a stochastic matrix
(Definition~\ref{def:tolmap}\,(iii)), and $c=1-\lmax(B)\in[0,1]$
(Definition~\ref{def:tolmap}\,(iv)). At $\varepsilon\ge1$ the corresponding
test is not merely loose, it is \emph{vacuous or degenerate}, and namably so:
Definition~\ref{def:tolmap}\,(i) would admit every Hermitian family with
$\lmin\ge-1$; \eqref{eq:eps2rank} would return $\rk_{\varepsilon_2}(A_k)=0$ for
\emph{every} port Gram, so the rank check could never refuse (the
under-reporting direction of Remark~\ref{rem:ranktwosided} taken to its
limit); (iii) would call every declared matrix monotone; (iv) would admit every
computed $c$; and (v) would admit a $\Gmat$ of operator norm $\sqrt2$.
$\varepsilon_{\mathrm{gap}}$ is bounded for the same reason at one remove: it
bounds a slack in $\pmc\in[0,1]$ (Theorem~\ref{thm:emitted-bound}), so a value
$\ge1$ permits the emitted bound to sit below the attainable one by the whole
range of the quantity, which is no constraint. The two members the exact
window adds are bounded on the same argument, up to a computed
$\widehat\Bnd^{\ast}$ that rounding puts outside $[0,1]$, which is not clamped
into that interval and is not claimed to lie in it
(Remark~\ref{rem:exacteig}). $\varepsilon_6$ bounds the positive part of
$\Bnd(n)-\widehat\Bnd^{\ast}(n)$, where $\Bnd(n)\le1$
(Theorem~\ref{thm:multiclick-bound}) and the exact value that
$\widehat\Bnd^{\ast}(n)$ evaluates lies in $[0,1]$
(Lemma~\ref{lem:exactbasic} with Definition~\ref{def:setting-family}); so at
$\varepsilon_6\ge1$ the test $r_6\le\varepsilon_6$ can refuse only an
evaluation that puts $\widehat\Bnd^{\ast}(n)$ below the evaluated $\Bnd(n)$
less $1$, which is at most $0$ up to rounding, outside the interval where the
exact value lies. $\varepsilon_{\mathrm{eig}}$ is subtracted from a quantity
whose exact value $\Bnd^{\ast}$ lies in $[0,1]$ at every receiver
(Lemma~\ref{lem:exactbasic} with \eqref{eq:mcpos}, which
Definition~\ref{def:setting-family} gives). So at
$\varepsilon_{\mathrm{eig}}\ge1$ the second entry of \eqref{eq:elldef} is
positive only where a computed $\widehat\Bnd^{\ast}$ exceeds $1$; elsewhere it
is never positive and the exact window is never the binding term: the
declaration pays the full cost of the computation and buys nothing, which is
degenerate in the same namable sense as the rest of this list.

Two arguments use the upper bound. Remark~\ref{rem:cdirection}'s even-$n$
estimate uses $\varepsilon_2^{\,n}\le\varepsilon_2^{\,2}$ for every even
$n\ge2$, which is exactly $\varepsilon_2\le1$; and the justification of
\eqref{eq:eps2rank} in Remark~\ref{rem:ranktwosided}, that the spectrum it
thresholds is normalized to $[0,1]$, is an argument about a threshold
\emph{inside} that interval.

The domain is a floor, not sizing advice. Nothing here says a tolerance near
$1$ is sensible; Remark~\ref{rem:ranktwosided} records that the choice of
$\varepsilon_2$ decides which declarations are accepted, and
Proposition~\ref{prop:dark-slack-bound} gives the inequality a declarer sizes $\varepsilon_2$
and $\varepsilon_4$ against. What the declared domain buys is that every error
term this contract states holds on every value the contract admits.
\end{remark}
